%
% 000-session11.tex
% Presentation Session 11: Practice 3 - Seq Server
%
% Compile to .pdf with LaTeX (pdflatex)
% Requires Beamer package (latex-beamer)
%

\documentclass[xcolor=table]{beamer}
\usetheme{Warsaw}
\beamertemplatenavigationsymbolsempty
\setbeamertemplate{headline}{}
\setbeamerfont{title}{size=\large}
\useoutertheme{infolines}
\usepackage[english]{babel}
\usepackage[utf8]{inputenc}
\usepackage{graphics}
\usepackage{amssymb}
\usepackage{multirow}
\usepackage{xcolor}
\usepackage{framed}
\usepackage{hyperref}

\definecolor{shadecolor}{RGB}{180,180,180}

%%%%%%%%%%%%%%%%%%%%%%%%%%%%%%%%%%%%%%%%%%%%%%%%%%%%%%%%%%%%%%%%
% Python Highlighting
%%%%%%%%%%%%%%%%%%%%%%%%%%%%%%%%%%%%%%%%%%%%%%%%%%%%%%%%%%%%%%%%
\DeclareFixedFont{\ttb}{T1}{txtt}{bx}{n}{10}
\DeclareFixedFont{\ttm}{T1}{txtt}{m}{n}{10}

% Custom colors
\usepackage{color}
\definecolor{deepblue}{rgb}{0,0,0.5}
\definecolor{deepred}{rgb}{0.6,0,0}
\definecolor{deepgreen}{rgb}{0,0.5,0}
\definecolor{grey1}{rgb}{0.5,0.5,0.5}

\usepackage{listings}

% Python style for highlighting
\newcommand\pythonstyle{\lstset{
language=Python,
basicstyle=\ttm\small,
morekeywords={self},
keywordstyle=\ttb\color{deepblue},
emphstyle=\ttb\color{deepred},
stringstyle=\color{deepgreen},
commentstyle=\small\color{grey1}\ttm,
frame=single,
showstringspaces=false
}}

% Python environment
\lstnewenvironment{python}[1][]
{
\pythonstyle
\lstset{#1}
}
{}

% Python for inline
\newcommand\pythoninline[1]{{\pythonstyle\lstinline!#1!}}

\newcommand{\asignatura}{Session 11: Practice 3 -- Seq Server}
\newcommand{\grado}{Programming in Network Environments}
\newcommand{\curso}{2025-2026}

\title[\asignatura]{\asignatura}
\subtitle{\grado}
\author{URJC}
\institute{GSyC}
\date{Course \curso}

\AtBeginSection[]
{
\begin{frame}<beamer>
\begin{center}
{\LARGE \insertsection}
\end{center}
\end{frame}
}

\begin{document}

\frame{
\maketitle
}

\begin{frame}
\frametitle{Contents}
\tableofcontents
\end{frame}

%%%%%%%%%%%%%%%%%%%%%%%%%%%%%%%%%%%%%%%%%%%%%%%%%%%%%%%%%%%%%%%%
%%%%%%%%%%%%%%%%%%%%%%%%%%%%%%%%%%%%%%%%%%%%%%%%%%%%%%%%%%%%%%%%
% SECTION 1: Session 10 Review
%%%%%%%%%%%%%%%%%%%%%%%%%%%%%%%%%%%%%%%%%%%%%%%%%%%%%%%%%%%%%%%%
%%%%%%%%%%%%%%%%%%%%%%%%%%%%%%%%%%%%%%%%%%%%%%%%%%%%%%%%%%%%%%%%

\section{Review: Session 10 -- Client-Server}

\begin{frame}
\frametitle{What We Learned in Session 10 (1/2)}

\textbf{Session 10: Introduction to Server Programming}:

\begin{itemize}
\item \textbf{Servers} use \textbf{two sockets}:
  \begin{itemize}
  \item A \textbf{listening socket} (ls) -- waits for incoming connections
  \item A \textbf{communication socket} (cs) -- created per client
  \end{itemize}
\vspace{0.15cm}
\item \textbf{Server setup steps}:
  \begin{enumerate}
  \item \texttt{socket()} -- Create the socket
  \item \texttt{bind(IP, PORT)} -- Assign address
  \item \texttt{listen()} -- Start listening
  \item \texttt{accept()} -- Wait for a client (blocks)
  \item \texttt{recv()} / \texttt{send()} -- Communicate
  \item \texttt{close()} -- Close client socket, loop back
  \end{enumerate}
\vspace{0.15cm}
\item \textbf{Address reuse}: \pythoninline{ls.setsockopt(socket.SOL_SOCKET, socket.SO_REUSEADDR, 1)}
\end{itemize}

\end{frame}

\begin{frame}[fragile]
\frametitle{What We Learned in Session 10 (2/2)}

\textbf{Happy Server} -- always responds with the same message:

\vspace{0.1cm}

\begin{python}
import socket

ls = socket.socket(socket.AF_INET,
                   socket.SOCK_STREAM)
ls.setsockopt(socket.SOL_SOCKET,
              socket.SO_REUSEADDR, 1)
ls.bind(("127.0.0.1", 8080))
ls.listen()

while True:
    (cs, client_ip_port) = ls.accept()
    msg = cs.recv(2048).decode()
    response = "HELLO. I am the Happy Server\n"
    cs.send(response.encode())
    cs.close()
\end{python}

\vspace{0.1cm}

\textbf{Exercises}: Echo server (1--3) + test client (4)

\end{frame}

%%%%%%%%%%%%%%%%%%%%%%%%%%%%%%%%%%%%%%%%%%%%%%%%%%%%%%%%%%%%%%%%
%%%%%%%%%%%%%%%%%%%%%%%%%%%%%%%%%%%%%%%%%%%%%%%%%%%%%%%%%%%%%%%%
% SECTION 2: Session 11 Overview
%%%%%%%%%%%%%%%%%%%%%%%%%%%%%%%%%%%%%%%%%%%%%%%%%%%%%%%%%%%%%%%%
%%%%%%%%%%%%%%%%%%%%%%%%%%%%%%%%%%%%%%%%%%%%%%%%%%%%%%%%%%%%%%%%

\section{Session 11: Practice 3 -- Seq Server}

\begin{frame}
\frametitle{Session 11 Goals}

\textbf{In this session we will}:

\begin{itemize}
\item Develop our own \textbf{Seq Server} to work with \textbf{DNA sequences}
\item Define a \textbf{communication protocol} with specific commands
\item Reuse \textbf{Seq Class} from Practice 1 and \textbf{Client Class} from Practice 2
\item Make DNA calculations \textbf{available over the network}
\end{itemize}

\vspace{0.4cm}

\textbf{Time}: 2 hours

\vspace{0.2cm}

\textbf{Files}: \texttt{P03/SeqServer.py} and \texttt{P03/test-client.py}

\vspace{0.2cm}

\textbf{Important}: Finish all S10 exercises before starting!

\end{frame}

\begin{frame}
\frametitle{The Seq Server Protocol}

\begingroup\small
\begin{center}
\begin{tabular}{|l|l|l|}
\hline
\textbf{Command} & \textbf{Argument} & \textbf{Response} \\
\hline
\texttt{PING} & (none) & ``OK'' \\
\hline
\texttt{GET n} & n: 0--4 & Sequence number n \\
\hline
\texttt{INFO seq} & A DNA sequence & Length + base counts (\%) \\
\hline
\texttt{COMP seq} & A DNA sequence & Complement sequence \\
\hline
\texttt{REV seq} & A DNA sequence & Reverse sequence \\
\hline
\texttt{GENE name} & Gene name & Full gene sequence \\
\hline
\end{tabular}
\end{center}
\endgroup

\vspace{0.2cm}

\textbf{Valid gene names} for GENE: \texttt{U5}, \texttt{ADA}, \texttt{FRAT1}, \texttt{FXN}, \texttt{RNU6\_269P}

\vspace{0.2cm}

\textbf{Testing}: Use \texttt{echo "PING" | nc 127.0.0.1 8080}

\vspace{0.2cm}

The server is built \textbf{incrementally} -- one command per exercise.

\end{frame}

\begin{frame}[fragile]
\frametitle{Exercises 1--2: PING and GET}

\textbf{Exercise 1: PING command}
\begin{itemize}
\item Parse the request message; if \texttt{PING} is detected, respond with \texttt{"OK!\textbackslash n"}
\item Print \texttt{"PING command!"} on the server console
\item Test: \texttt{echo "PING" | nc 127.0.0.1 8080}
\end{itemize}

\vspace{0.3cm}

\textbf{Exercise 2: GET command}
\begin{itemize}
\item \texttt{GET n} returns the sequence at index \texttt{n} (0--4)
\item Store 5 test sequences in a \textbf{list}
\item Parse the command and extract the number argument
\item Test: \texttt{echo "GET 2" | nc 127.0.0.1 8080}
\end{itemize}

\vspace{0.2cm}

\textbf{Hint -- parsing the command}:

\begin{python}
cmd = msg.strip().split(" ", 1)
command = cmd[0]  # e.g. "GET"
if len(cmd) > 1:
    argument = cmd[1]  # e.g. "2"
\end{python}

\end{frame}

\begin{frame}[fragile]
\frametitle{Exercises 3--4: INFO and COMP}

\textbf{Exercise 3: INFO command}
\begin{itemize}
\item \texttt{INFO seq} returns information about the sequence
\item Use the \textbf{Seq Class} to compute the data
\item Response format:
\end{itemize}

\begin{verbatim}
  Sequence: AACCGTA
  Total length: 7
  A: 3 (42.9%)
  C: 2 (28.6%)
  G: 1 (14.3%)
  T: 1 (14.3%)
\end{verbatim}

\vspace{0.2cm}

\textbf{Exercise 4: COMP command}
\begin{itemize}
\item \texttt{COMP seq} returns the \textbf{complement} of the sequence
\item Use the \textbf{Seq Class} to calculate it
\item Test: \texttt{echo "COMP AACCGTA" | nc 127.0.0.1 8080}
\end{itemize}

\end{frame}

\begin{frame}
\frametitle{Exercises 5--6: REV and GENE}

\textbf{Exercise 5: REV command}
\begin{itemize}
\item \texttt{REV seq} returns the \textbf{reverse} of the sequence
\item Use the \textbf{Seq Class} to calculate it
\item Test: \texttt{echo "REV AACCGTA" | nc 127.0.0.1 8080}
\end{itemize}

\vspace{0.4cm}

\textbf{Exercise 6: GENE command}
\begin{itemize}
\item \texttt{GENE name} returns the \textbf{full sequence} of a gene
\item Valid names: \texttt{U5}, \texttt{ADA}, \texttt{FRAT1}, \texttt{FXN}, \texttt{RNU6\_269P}
\item Use the \textbf{Seq Class} to read the gene from its file
\item Test: \texttt{echo "GENE U5" | nc 127.0.0.1 8080}
\end{itemize}

\end{frame}

\begin{frame}[fragile]
\frametitle{Exercise 7: Test Client}

Write a \textbf{client program} (\texttt{P03/test-client.py}) that tests \textbf{all commands}:

\begin{itemize}
\item Use the \textbf{Client Class} from P02 and the \texttt{talk()} method
\item For \texttt{INFO}, \texttt{COMP}, and \texttt{REV}, use the sequence from \texttt{GET 0}
\item Print the command being tested and the server response
\end{itemize}

\vspace{0.2cm}

\begin{python}
from Client0 import Client

c = Client("127.0.0.1", 8080)
print(c)
print("* Testing PING...")
print(c.talk("PING"))
print("* Testing GET...")
for i in range(5):
    print(f"GET {i}: {c.talk(f'GET {i}')}")
\end{python}

\end{frame}

\begin{frame}
\frametitle{Deliverables and Checklist}

\textbf{P03 folder should contain}:

\begin{itemize}
\item \textbf{SeqServer.py} -- The Seq Server with all 6 commands
\item \textbf{test-client.py} -- Client that tests all commands
\end{itemize}

\vspace{0.3cm}

\textbf{Checklist}:
\begin{itemize}
\item[$\square$] All S10 items are completed
\item[$\square$] SeqServer.py implements: PING, GET, INFO, COMP, REV, GENE
\item[$\square$] test-client.py tests all commands and prints results
\item[$\square$] Server uses the \textbf{Seq Class} from P01
\item[$\square$] Client uses the \textbf{Client Class} from P02
\item[$\square$] All files pushed to your remote Github repo
\end{itemize}

\vspace{0.3cm}

\textbf{Remember}: You are \textbf{building on previous work} (P01 + P02)!

\end{frame}

%%%%%%%%%%%%%%%%%%%%%%%%%%%%%%%%%%%%%%%%%%%%%%%%%%%%%%%%%%%%%%%%
%%%%%%%%%%%%%%%%%%%%%%%%%%%%%%%%%%%%%%%%%%%%%%%%%%%%%%%%%%%%%%%%
% SECTION 3: End
%%%%%%%%%%%%%%%%%%%%%%%%%%%%%%%%%%%%%%%%%%%%%%%%%%%%%%%%%%%%%%%%
%%%%%%%%%%%%%%%%%%%%%%%%%%%%%%%%%%%%%%%%%%%%%%%%%%%%%%%%%%%%%%%%

\begin{frame}
\begin{center}
\Huge{Questions?}

\vspace{1cm}

\Large{Time to code!}
\end{center}
\end{frame}

\end{document}
